% T29 source: Примеры базовой модели.tex lines 398--488
\detailedexampleheading{Variant 1. Interactive Work in a Single Stream}

The developer completes the tasks sequentially using an AI assistant. For each
task, the developer formulates a prompt, receives the result, reviews it, and
corrects it if necessary. There is no parallelism
($P = 1$). The configuration is
$\sigma_1 = (P{=}1,\ \text{interactive mode})$.

\paragraph{Parameters.}
\begin{table}[ht!]
\centering
\caption{Parameters of Variant 1}
\label{tab:practical-v1-params}
\begin{tabular}{clccc}
\toprule
$i$ & Task & $Z_i$ & $k_i$ & $h_i$ \\
\midrule
1 & API contract and request validation & 2 & 0.80 & 0.80 \\
2 & Persistence layer & 3 & 0.90 & 0.90 \\
3 & Business logic & 4 & 0.85 & 0.85 \\
4 & Background worker & 3 & 0.95 & 0.95 \\
5 & Observability and configuration & 2 & 0.75 & 0.75 \\
6 & Tests & 3 & 0.80 & 0.80 \\
7 & Packaging and deployment & 2 & 0.90 & 0.90 \\
\bottomrule
\end{tabular}
\end{table}

The coefficients $k_i$ encode moderate acceleration: AI assists in generating
boilerplate code and tests, but every result undergoes manual review. For tasks
with high architectural complexity (business logic and the background worker),
the coefficient is closer to one. The human component satisfies $h_i = k_i$:
in interactive mode, the developer retains the task context throughout its
execution (the agent component is $a_i = 0$).

\paragraph{Calculations.}
\begin{align*}
  \sum_{i=1}^{7} Z_i k_i &= 2{\cdot}0.80 + 3{\cdot}0.90 + 4{\cdot}0.85 + 3{\cdot}0.95 + 2{\cdot}0.75 + 3{\cdot}0.80 + 2{\cdot}0.90 \\
  &= 1.60 + 2.70 + 3.40 + 2.85 + 1.50 + 2.40 + 1.80 = 16.25.
\end{align*}

The weighted-average normalized duration is
\[
  \bar{k}_w = \frac{16.25}{19} \approx 0.855.
\]

The task weights in the graph are $w_i = Z_i k_i X$:
\[
  \{w_i\}_{i=1}^{7} = \{6.4,\; 10.8,\; 13.6,\; 11.4,\; 6.0,\; 9.6,\; 7.2\}~\text{h}.
\]

The total work is
\[
  W = \sum_{i=1}^{7} w_i = 16.25 \cdot 4 = 65.0~\text{h}.
\]

The idealized completion time for independent tasks, that is, the lower bound
based on total work when $P = 1$, is
\[
  T_{ai}^{(0)} = \frac{W}{P} = \frac{65.0}{1} = 65.0~\text{h}.
\]

The idealized speedup without accounting for dependencies is
\[
  S_E = \frac{76}{65.0} \approx 1.17.
\]

The critical path through the graph comprises Tasks
$1 \to 2 \to 3 \to 4 \to 6$:
\[
  L = 6.4 + 10.8 + 13.6 + 11.4 + 9.6 = 51.8~\text{h}.
\]

The lower bound on the makespan accounting for the dependency graph is
\[
  B_4 = \max\left\{\frac{W}{P},\; L,\; H\right\} = \max\{65.0,\; 51.8,\;65.0\} = 65.0~\text{h}.
\]

When $P=1$, the schedule is sequential. For example, the tasks may be
completed in the topological order $1,2,3,4,5,6,7$. The makespan equals the
total work:
\[
  T(\pi_1)=T^{*}=B_4=65.0~\text{h}.
\]

The achievable speedup accounting for the graph is
\[
  S_E^{\text{ach}} = \frac{76}{65.0} \approx 1.17.
\]

At level M4, the human-resource calculation is as follows because $h_i = k_i$:
\[
  H = X \sum_{i=1}^{7} Z_i h_i = 65.0~\text{h} = W,
  \qquad
  \bar{h}_w = \bar{k}_w \approx 0.855,
  \qquad
  \frac{1}{\bar{h}_w} \approx 1.17.
\]

\paragraph{Interpretation.}
Interactive mode provides direct control over the result. AI reduces the
effort required for each task ($\bar{k}_w \approx 0.855$), but the absence of
parallelism limits the overall effect. The critical path is not the active
constraint: when $P=1$, $W/P=65.0$~h exceeds $L=51.8$~h. The speedup is
$1.17\times$, corresponding to a savings of approximately 11 hours relative
to the 76-hour no-AI baseline.

At level M4, the achieved speedup of $1.17\times$ exactly equals the
human-resource ceiling $1/\bar{h}_w$. In interactive mode, the makespan equals
the total human time $H$. Consequently, further process reorganization cannot
yield an improvement without reducing $h$, that is, without transferring part
of the work to agents for autonomous execution.
